\section{Conclusion}

% We introduced SoundnessBench to evaluate a core capability for autonomous science: deciding whether a hypothesis and experimental design are scientifically sound before running expensive experiments. Across the models we tested, judgments are often sensitive to prompt framing: under standard prompting many models over-approve weak proposals, while stricter prompting can push some models toward broad over-rejection. At this early gate in the research pipeline, such instability can misallocate compute and attention. We view SoundnessBench as a step toward safer and more useful AI research agents, and as a tool for diagnosing where model calibration breaks down. The observed optimism tendency and prompt sensitivity suggest that reliable scientific judgment will likely require targeted training or calibration approaches beyond prompting alone.
\fh{We introduced SoundnessBench to evaluate a core capability for autonomous ML research agents: deciding whether a hypothesis and experimental design appear methodologically sound before running expensive experiments. Across the models we tested, judgments are often sensitive to prompt framing: under standard prompting many models over-approve weak proposals, while stricter prompting can push some models toward broad over-rejection. Additional controls for public-corpus contamination, identifier recognition, surface features, and slice-level concentration do not remove this pattern. At this early gate in the research pipeline, such instability can misallocate compute and attention. We view SoundnessBench as a step toward safer and more useful AI research agents, and as a tool for diagnosing where model calibration breaks down. The observed optimism tendency and prompt sensitivity suggest that reliable proposal-stage scientific judgment will likely require targeted training, calibration, or human-in-the-loop review beyond prompting alone.}

% \textbf{Limitations.} Our ground truth relies on reviewer soundness scores, and the benchmark covers a bounded slice of ML papers. Expanding to broader venues, modalities (e.g., code and logs), and longitudinal settings is an important next step. 
\fh{\textbf{Limitations.} Our ground truth relies on reviewer soundness sub-scores, which are expert signals but still imperfect proxies for proposal-only methodological validity because reviewers saw full papers, results, presentation quality, and framing. SoundnessBench should therefore be interpreted as measuring recoverable pre-execution soundness signals rather than exact full-paper review prediction or definitive post-execution research quality. The benchmark also covers a bounded slice of ML research drawn from ICLR; extending to other venues and scientific domains such as biology, chemistry, and social science is important before making claims about scientific soundness in general. Because the source corpus is public, perfect contamination control is impossible, although our ICLR 2026-only and identifier-removal analyses reduce this concern. Finally, our human audit is preliminary and does not establish a full expert-human ceiling. Expanding expert re-annotation, broader human audits, private or continuously refreshed test sets, richer modalities (e.g., code and logs), and longitudinal proposal-to-execution studies are important next steps.}
% In addition, scale and instruction-tuning ablations are limited to one model family (Qwen3.5), so cross-architecture generalization remains open.
